\section{Masks and initial numerator prescriptions}
\label{sec:prescriptions}

The centre assignments of Section~\ref{sec:centres} determine which
rows are available. We now remove a negligible set of their deficit
masks and define the weights at their first projections. The hard
weights have a uniform lower mean comparable to the totient density;
the additional soft conditions have arbitrarily small relative
first-mass cost. All prescriptions here are fixed before the
adaptive tables are constructed.

\subsection{Good deficit masks}

For a centre $L=L_i$ of band $B=B_i$, the assigned masks are
$d=L/v\in\mathcal F_i$. Probabilities below refer to the unrestricted
deficit law $\pi_L$, not to a uniform distribution on assigned rows.
Recall that its prime coordinates are independent and
$\Pr(\nu_p(d)\ge j)=p^{-j}$ for $j\le\nu_p(L)$.

\begin{definition}[Good masks]
\label{pre:good-masks}
An assigned mask $d$ is \emph{good} if it satisfies all applicable
conditions below.
\begin{enumerate}
\item Its excess prime-power part obeys
\[
 \prod_p p^{(\nu_p(d)-1)_+}\le e^{C_{\rm mask}B}.
\]
\item On every initial segment of the prime divisors of $L$ whose
reciprocal sum is at most $s_{\rm low}B$, the part of $d$ supported
on that segment has divisor count at most $e^{s'B}$.
\item Its large-prime harmonic cost satisfies
\[
 \sum_{\substack{p\mid d\\p>d_0B}}\frac{(\log p)^{.3}}p\le1.
\]
\item On high type, the last dense bin contains at least
$c\kappa Y_d$ prime divisors of $d$, where $c>0$ is a sufficiently
small fixed absolute constant.
\end{enumerate}
\end{definition}

An \emph{omitted} prime means a prime dividing the deficit mask:
one copy is absent from the row denominator, not necessarily the
entire prime power of the centre. In all tagged or tail ranges below,
good masks have deficit at most one at each prime. Indeed these
primes exceed $e^{JB}$, or are still larger in terms of their starting
height, and $J>C_{\rm mask}$.

\begin{lemma}[Negligible mask deletion]
\label{pre:mask-purge}
With the order of choices in Section~\ref{cen:hierarchy}, deleting
all assigned masks that are not good loses $o(W)$ as $k_0\to0$.
Uniformly for each retained centre, its unrestricted failure
probability is $o(\delta_L)$.
\end{lemma}

\begin{proof}
For the first condition take a fixed small positive moment of the
excess prime-power part. The geometric deficit tails give a bounded
Euler product: the local contribution beyond one starts at exponent
two and is $1+O(p^{-2+\epsilon})$. With this fixed moment exponent,
choose $\epsilon C_{\rm mask}>2$. Markov's inequality then gives
failure probability $\ll e^{-\epsilon C_{\rm mask}B}$.

For the second condition it is enough to test the largest allowed
initial segment, or the entire support if its reciprocal sum is
small enough. Every fixed moment of the divisor count on this
segment is at most
\[
 \exp\left(C_{\rm moment}\sum_{p\text{ in the segment}}\frac1p\right).
\]
Choose a fixed moment $M$ with $Ms'>3$, then choose $s_{\rm low}$
so small that $Ms'-C_Ms_{\rm low}>2$. Markov's inequality makes
the failure probability at most $e^{-2B}$. This moment choice
occurs between the choices of $s'$ and $s_{\rm low}$ and does not
depend on any later tag or model parameter.

For the third condition the logarithm of the moment generating
function with coefficient $2B$ is at most
\[
 \sum_{p>d_0B}
 \frac{\exp(2B(\log p)^{.3}/p)-1}{p}=o(B).
\]
To see the last bound, split at $B^2$. Below that point one can
bound even the integer sum by
$O(\log B\exp(C_{d_0}(\log B)^{.3}))=o(B)$. Above it the
exponent is small and the integer sum is
$O(B\sum_{n>B^2}(\log n)^{.3}/n^2)=o(B)$.
Thus the failure probability is $e^{-2B+o(B)}$. Summing over all
integers has made the estimate uniform over every centre prime
support: no factor depending on $\log\log L$ occurs. The fixed
$d_0$ affects only how large $B$ must be for the $o(B)$ estimate,
which is permitted because $k_0$ is chosen last.

For the last condition, independence and the reciprocal sum in
the dense bin give
\[
 \mathbb E\exp\left(-\sum_{\substack{p\mid L\\p\in\mathcal P(Y_d)}}
 \mathbf1_{\{p\mid d\}}\right)
 =\prod_{\substack{p\mid L\\p\in\mathcal P(Y_d)}}
 \left(1-\frac{1-e^{-1}}p\right)
 \le e^{-c_1\kappa Y_d}.
\]
For sufficiently small $c$, failure has probability
$e^{-c_2\kappa Y_d}$. Since $Y_d>D_0B$ on high type, choose
$D_0$ after $\kappa$ so that $c_2\kappa D_0>2$.

Take $C_{\rm mask}$ large in the first estimate. All four
probabilities are then $o(e^{-.64B})$, uniformly as the minimum
band tends to infinity. Since $B_i^*<2B_i$,
\eqref{eq:centre-lower} gives $\delta_L\ge e^{-.64B}$.
Multiplying each failure probability by $k_L$, and using
$m_L=k_L\delta_L$, proves the claimed relative and total losses.
\end{proof}

From now on assigned masks are good. Nonnegative upper bounds may
still extend sums to the unrestricted mask laws.

\subsection{The first projection and its currency}

At centre $L_i$, let $P_I$ be the product of the full prime powers
of $L_i$ with $\log\log p\ge Y_0(i)$. Define the tagged factor
$P_M$ as follows. On high type it contains precisely the full prime
powers in the last dense bin. On medium type it contains all full
prime powers at $p>e^{JB}$ below the tail $P_I$. For medium type
partition the tagged primes into the ordinary bins
$\ell B\le Y'<Y_0(i)$ and the extra bin
\[
 e^{JB}<p<\exp(\exp(\ell B)).
\]
Empty bins require no action. Put $P=P_MP_I$ and write $L_i=S_iP$;
these factors have disjoint prime supports.

The first projection restores only factors omitted in the unprocessed
tail; the low and tagged deficits remain in its denominator.
For an assigned row, decompose its denominator and its deficit into
\[
 v=V(P_M/m)(P_I/a),\qquad
 V\mid S_i,\quad m\mid P_M,\quad a\mid P_I.
\]
Its \emph{birth projection} is the grid
\begin{equation}
 \begin{gathered}
 x=\frac{z+a\gamma}{VP/m},\quad z\in\mathbb Z,
 \\
 r_i\pi_{P_I}(a)(VP/m)R(VP_M/m)=r_i\phi(v).
 \end{gathered}
 \label{pre:birth-grid}
\end{equation}
The true row points are exactly the subarray $a\mid z$. The identity
in \eqref{pre:birth-grid} follows from multiplicativity across the
three disjoint prime supports. It is the first-mass normalization:
the entire projection receives the single coefficient
$\pi_{P_I}(a)$, while the baseline numerator density is
$R(VP_M/m)$. The coefficient is not an additional numerator test.

\subsection{Tagged hard factors}

Write
\[
 \alpha=a\gamma,\qquad d^0=(S_i/V)m,\qquad
 \mathcal L_i=\lfloor d^0\alpha\rfloor.
\]
These quantities depend on the current row and its masks, not just
on the centre whose index appears in $\mathcal L_i$.
For a rational whose denominator is a unit modulo a prime, its
congruence class is interpreted by inversion of that denominator.

\begin{construction}[Tagged numerator weights]
\label{pre:tag-prescription}
In each tag bin in which $m$ has at least one prime divisor, sample
an anchor uniformly from those divisors. In an ordinary bin of
height $Y'$, try to find a reduced rational $F_i$ satisfying
\[
 F_i\equiv\mathcal L_i\pmod {p_{\rm anchor}},\qquad
 |F_i|,\ \operatorname{den}(F_i)
 \le\exp(\exp(h_HY')).
\]
In the extra medium bin, replace this bound by $e^{EB}$.
The key $F_i$ is local to this bin and anchor choice; different bins
may use different keys despite the suppressed bin index.
If such a rational is found, then at every full tag factor in that
bin, namely $p\mid P_M$ with $p\nmid m$, forbid
\[
 d^0z+\mathcal L_i\equiv F_i\pmod p.
\]
If reconstruction fails, or if the bin has no anchor because its
tag mask is empty, use the scalar weight $R_p$ at those full
factors. At every omitted tag factor with positive residual
exponent, use the scalar $R_p$ as well. Average jointly over the
independent anchor choices in the different bins.
\end{construction}

There is at most one reconstructed rational. Write the height bound
in its bin as $Q$. If $F_1=n_1/q_1$ and $F_2=n_2/q_2$ were two
candidates, then $q_i\le Q$ and $|n_i|\le Qq_i$, so
\[
 |n_1q_2-n_2q_1|\le 2Qq_1q_2\le2Q^3.
\]
The anchor prime exceeds $2Q^3$: this follows from $h_H<1$ in
ordinary bins and $J>4E$ in the extra bin, once the bands are large.
It also exceeds $Q$, so the candidate denominators are units.
The congruence forces the displayed integer to vanish, proving
uniqueness. The same size comparison applies to every prime where
the candidate is used. At a full tag factor $d^0$ is also a unit,
so a successful reconstruction forbids exactly one numerator
residue. Anchor averaging makes the final prescription
deterministic.

\subsection{Low rational models}

The numerator restrictions in this subsection use a rational model of
the fixed projected phase. Shift-adapted primitivity also occurs in
Chow and Technau~\cite[Definition~1.20]{ChowTechnau2024}; here the
models, thresholds and gcd bounds are defined for the present rows
and their projections.

For the projected numerator width and relative rational height put
\[
 \psi=(VP/m)r_i,\qquad
 H_v(\zeta)=2+q+\frac{|\alpha-\zeta|}{\psi},
 \quad \zeta=n_0/q\text{ reduced},\quad q\ge1.
\]
Here the scalar $\psi$ is the projected numerator width, not the
original tolerance function in Theorem~\ref{thm:main}.
The subscript $v$ includes the row's fixed width and birth
projection. In particular $\psi\le k_L/m$.

\begin{definition}[Recognized models]
\label{pre:recognized-models}
The two medium broad tiers are alternatives; the narrow tier is an
additional test. High type uses only its single broad tier.
On medium type, use the large broad tier if
$\psi\le e^{-4EB}$, recognizing a rational when
$H_v(\zeta)\le e^{EB}$. Otherwise use the small broad tier only if
$\psi\le e^{-4h_0B}$, with threshold $H_v(\zeta)\le e^{h_0B}$.
Also test the narrow tier
\[
 q\le e^{EB},\qquad |\alpha-\zeta|\le\psi e^{-6EB}.
\]
On high type use just one broad tier,
\[
 H_v(\zeta)\le\exp(\exp(h_HY_d)).
\]
\end{definition}

\begin{lemma}[Uniqueness and separation of models]
\label{pre:model-uniqueness}
Each tier recognizes at most one rational. The same statement
holds across rows with a fixed $\alpha$ and the same tier
parameters: a fixed $B$ on medium type, or a fixed $Y_d$ on high
type, even if the latter rows have different bands. At fixed
$\alpha,B$, a large broad model and a narrow model recognized
on different medium rows must also coincide. If a medium row has
distinct broad and narrow models,
the broad tier is the small one and the narrow model has
denominator at least $e^{cB}$ for a fixed $c>0$.
\end{lemma}

\begin{proof}
Distinct reduced rationals of denominators at most $Q$ differ by
at least $Q^{-2}$. In a medium broad tier with exponent
$h\in\{E,h_0\}$, the error from $\alpha$ is at most $e^{-3hB}$
and the denominator is at most $e^{hB}$, proving uniqueness.
In the narrow tier, the centre scale bound gives
\[
 \psi\le k_L\le e^{.2B+O(1)},
\]
so its error is at most $e^{-(6E-.2)B+O(1)}$, again smaller than
half the rational separation. This estimate is uniform across
rows at that fixed band.

For high type, the good-mask occupancy condition supplies a tag
prime in $m$, and hence
\[
 \psi\le\exp(.2B+O(1)-\exp(Y_d)).
\]
For all such rows at fixed $Y_d$, put
$Q=\exp(\exp(h_HY_d))$. Their model errors are uniformly at most
\[
 \exp\bigl(.2Y_d/D_0+O(1)-\exp(Y_d)+\exp(h_HY_d)\bigr)
 <\tfrac12 Q^{-2},
\]
for sufficiently large $Y_d$. This proves the asserted cross-row
uniqueness because $h_H<1$.

The large medium broad model and narrow model cannot be distinct,
even on different rows with the same $\alpha,B$: their two error
bounds sum to less than $e^{-2EB}$, and both denominators are at
most $e^{EB}$. If the small broad model has denominator $q_b$
and a distinct narrow model has denominator $q_n$, then
\[
 \frac1{q_bq_n}
 \le |\zeta_b-\zeta_n|
 \le e^{-3h_0B}+e^{-(6E-.2)B+O(1)}.
\]
Using $q_b\le e^{h_0B}$ gives $q_n\gg e^{2h_0B}$, which is
the last assertion.
\end{proof}

For every recognized model $n_0/q$, require at each $p\mid V$
\begin{equation}
 p\nmid qz+n_0
 \quad\text{if}\quad p>H_v(\zeta)^{\epsilon_s}
 \text{ and }p\nmid q.
 \label{pre:low-hard-tests}
\end{equation}
If no exclusion is active at $p$, multiply by $R_p$ instead. When
two models impose exclusions, impose both. The resulting low
factors, together with Construction~\ref{pre:tag-prescription},
are called the \emph{hard birth weights}.

\begin{lemma}[Hard first-mass bounds]
\label{pre:hard-mean}
There is $c_{\rm hard}>0$, independent of $k_0$ once it is
sufficiently small, such that the mean hard birth weight on a row's
projected numerator grid lies in
\[
 [c_{\rm hard}R(VP_M/m),\ R(VP_M/m)].
\]
Consequently its first mass, including the single coefficient
$\pi_{P_I}(a)$ and the spatial width, lies between
$c_{\rm hard}\phi(v)r(v)$ and $\phi(v)r(v)$.
\end{lemma}

\begin{proof}
Fix the anchor choices. At a prime with one active exclusion the
uniform residue mean is exactly $R_p$; the same is true of its
scalar replacement. At an omitted tag factor the scalar is
present exactly when that prime remains in the projected
non-tail denominator. The prime coordinates of the numerator
are uniform and independent by CRT.

Only two distinct medium models can reduce the mean below this
one-factor product. By Lemma~\ref{pre:model-uniqueness}, their
additional narrow-model hard primes exceed $e^{c\epsilon_sB}$.
On medium type every prime of $V$ is at most $e^{JB}$. Thus the
reciprocal sum of these additional primes is bounded by a
constant independent of $B$ and $k_0$. They are also all large
when $k_0$ is small; in particular two distinct active exclusions
cannot occur at $p=2$. The local mean with two exclusions is at
least $1-2/p$, whose ratio to $R_p$ is $1-1/(p-1)$.
Multiplying these ratios gives a positive uniform lower constant.
The upper mean is always at most the one-factor product.
These bounds persist after averaging the anchor choices.
The last assertion is \eqref{pre:birth-grid}.
\end{proof}

\subsection{Soft tests and the remaining initial deletion}

In addition to the hard weights, require for every recognized model
\begin{equation}
 (qz+n_0,V)\le H_v(\zeta)^{h_s}.
 \label{pre:soft-tests}
\end{equation}

\begin{lemma}[Relative cost of the soft tests]
\label{pre:soft-loss}
Under the probability measure obtained by normalizing a row's
hard birth weights, the probability that at least one of the tests
\eqref{pre:soft-tests} fails is $O(\epsilon_s/h_s)$, with an absolute
implied constant. In particular their relative first-mass loss can
be made arbitrarily small compared with $\mu$ by the stated
choice of $\epsilon_s$.
\end{lemma}

\begin{proof}
Fix a recognized model $n_0/q$ and the static anchor choices.
For $p\mid q$, reducedness gives $p\nmid n_0$, so that prime does
not divide $qz+n_0$. For $p\nmid q$ above the threshold
$H_v^{\epsilon_s}$, the model's own hard exclusion already makes
its valuation zero. At a remaining prime, at most the other
model can impose a nonconstant exclusion. Conditioning on it
inflates each geometric valuation tail by at most an absolute
factor. The scalar factors and the tag factors do not change
this conclusion. Therefore
\[
 \mathbb E_{\mathrm{hard}}\log(qz+n_0,V)
 \ll\sum_{p\le H_v^{\epsilon_s}}\frac{\log p}{p-1}
\ll\epsilon_s\log H_v.
\]
For the last inequality, if $H_v^{\epsilon_s}<2$ the prime sum
is empty; otherwise $\log(2H_v^{\epsilon_s})\ll
\epsilon_s\log H_v$. At $p=2$ conditioning on the one possible
other exclusion costs at most a factor two, so the same absolute
tail bound applies. If $n_0=0$, reducedness gives $q=1$; the
valuation tails are still the same geometric tails, truncated at
the exponents in $V$, including the numerator residue zero.
Markov's inequality gives failure probability
$O(\epsilon_s/h_s)$ for this model. There are at most two
recognized models, so the union bound proves the assertion.
For each fixed anchor choice this bounds failed hard mass by
$O(\epsilon_s/h_s)$ times its hard mass. Integrate this inequality
over anchor choices before normalizing. The normalized average
is a mixture of the conditioned laws, and has the same relative
bound, without an additional factor $c_{\rm hard}^{-1}$.
\end{proof}

One additional medium-type birth deletion is needed: it removes
small neighborhoods of rational points after specified integer
dilations. Its explicit definition and cost proof are given in
Section~\ref{sec:direct}. It uses only the aggregate lists of good
rows and squarefree moduli of the low numerator. In particular it
is specified before, and independently of, all running-table
deletions or estimates.

All birth requirements are multiplied on each entry, with weight
zero when a required exclusion or test fails; the anchor average
is joint over that entry's prescriptions. No static test uses a
tail-prime coordinate of the projected numerator $z$. Thus the
birth arrays retain all unprocessed-prime translation symmetries,
which are the input to the table construction in
Section~\ref{sec:tables}.
